\section{Methodology}
\label{sec:method}

We conduct four experiments that progressively isolate the contributions of pretraining/SFT and RLHF to verbosity and blandness. \Secref{sec:data} describes our datasets, \secref{sec:prism_analysis} analyzes real human preference diversity, \secref{sec:dpo_training} details our DPO training procedure, and \secref{sec:eval} describes our evaluation protocol.

\subsection{Datasets}
\label{sec:data}

\para{\prism.} We use the \prism pluralistic alignment dataset \citep{kirk2024prism}, which contains 68,371 utterances from 1,396 unique users. Each conversation turn includes two model responses scored by the user, with explicit chosen/rejected labels. We retain 1,393 users with at least three preference pairs for our per-user analysis.

\para{\personalllm.} For DPO training, we use \personalllm \citep{personalllm2024}, which provides 9,402 training prompts and 1,000 test prompts, each with 8 model responses scored by 10 different reward models that simulate annotators with diverse preference profiles. The 10 reward models exhibit varying length-score correlations: \gemma has $r{=}0.77$ (strong length preference), \beaver has $r{=}-0.16$ (slight brevity preference), and \oasst has $r{=}0.51$ (moderate length preference).

\subsection{Experiment 1: Per-User Preference Analysis}
\label{sec:prism_analysis}

For each \prism user, we compute: (1) the mean word count of chosen and rejected responses, (2) the length preference ratio (chosen length / rejected length), and (3) the Spearman correlation between response scores and word counts. We test whether individual users differ significantly in length preferences using the Kruskal-Wallis $H$ test.

\subsection{Experiment 2: DPO Training}
\label{sec:dpo_training}

We select three \personalllm annotators with maximally diverse length preferences and construct DPO training pairs by selecting the highest- and lowest-scored responses for each prompt according to each annotator's reward model. We also create an aggregate condition by averaging scores across all 10 annotators. Each condition uses 800 training pairs.

\para{Model and training.} We use \qwen as the base model, which is already instruction-tuned (SFT'd). We apply LoRA \citep{hu2022lora} with rank 16, $\alpha{=}32$, targeting the query and value projection layers. DPO uses $\beta{=}0.1$ as the KL penalty strength. We train for a single epoch with learning rate $5{\times}10^{-5}$, batch size 4 with 4 gradient accumulation steps, and maximum sequence length of 768 tokens. All training uses a single NVIDIA RTX A6000 GPU (49GB). Training takes approximately 5 minutes per model.

\para{Conditions.} We train four DPO models plus the base model as a no-RLHF control:

\begin{table}[h]
    \centering
    \caption{Experimental conditions. Length preference indicates the Spearman correlation between the annotator's scores and response word counts.}
    \begin{tabular}{@{}llr@{}}
        \toprule
        \textbf{Condition} & \textbf{Description} & \textbf{Length Pref.\ ($r$)} \\
        \midrule
        Base & \qwen, no DPO & --- \\
        \gemma & DPO on strong length-loving annotator & 0.77 \\
        \beaver & DPO on brevity-preferring annotator & $-$0.16 \\
        \oasst & DPO on moderate annotator & 0.51 \\
        Aggregate & DPO on mean of all 10 annotators & ${\approx}0.55$ \\
        \bottomrule
    \end{tabular}
    \label{tab:conditions}
\end{table}

\subsection{Experiments 3 and 4: Evaluation}
\label{sec:eval}

\para{Generation.} We generate responses on 100 held-out test prompts from \personalllm using temperature 0.7 and top-$p$ 0.9, yielding 500 total responses across 5 models.

\para{Automatic metrics.} We measure: (1) mean response length in words, (2) distinct-$n$ for $n \in \{1, 2\}$ (ratio of unique $n$-grams to total), (3) self-BLEU (average BLEU between pairs of outputs from the same model), (4) type-token ratio (TTR), and (5) inter-model TF-IDF cosine similarity.

\para{LLM judge.} We use \gptjudge to evaluate 50 randomly sampled responses per model on four dimensions (1--5 scale): helpfulness, verbosity, distinctiveness, and overall quality.

\para{Statistical tests.} We use Mann-Whitney $U$ tests to compare output length distributions between conditions, with Cohen's $d$ as the effect size measure. We set $\alpha = 0.05$.
